% TODO
% Response to referee 1, comment 7
Plan:
\begin{enumerate}
\item This is easy to refute, we just need Ajesh to write some pretty equations and show that this contribution goes like $(\delta|P|)^2$.
\end{enumerate}

We thank the referee for considering this, this is a point which was not stated clearly in the manuscript.
The crucial point is that the amplitudes of the collective oscillations are small compared to the equilibrium values.
Working to linear order, we show in the Supplemental Material that, indeed, the anisotropy-gapped pseudo-Goldstone mode does not couple to a longitudinal polarization fluctuation $\delta P_y$.
This has been noted in Ref.~\cite{} as well.
What happens at quadratic order?
Let us assume that the polarization follows the spin-ordering via: $\vec{P}_i \propto \hat{x} \times (\vec{S}_i \times \vec{S}_{i+1})$, which holds if the optical phonon frequency is much larger than the electromagnon frequencies (which is true in the present case).
When the pseudo-Goldstone magnon is excited, $\vec{S}_i = S_0 \left( \cos \varphi_i \hat{x} + \sin \varphi_i \hat{y} + \theta \sin \varphi_i \hat{z}\right)/\sqrt{1 + \theta^2 \sin^2 \varphi_i} $, where $\varphi_i \equiv \vec{Q} \cdot \vec{R}_i$.
Averaging over all sites, we obtain the induced polarizations: $\delta P_z = \theta P_0$ and $\delta P_y \propto -\frac{1}{2} \theta^2 P_0$. 
For $\theta(t) = \Theta\cos(\omega_{\mathrm{EM}} t)$, the longitudinal component therefore oscillates at $2\omega_{\mathrm{EM}}$, in agreement with the referee's statement.
However, this occurs at quadratic order in the small fluctuation $\theta$, which we expect makes the corresponding signal small in the experiment.

Importantly, we note that this does \textcolor{red}{not} affect our arguments for the identification of the Higgs mode, since the Higgs-mode frequency is much smaller than twice the pseudo-Goldstone mode frequency.

\begin{changesmade}
We have emphasized this point in the main text and added details to the Supplementary Material of the revised manuscript.
\end{changesmade}

\textcolor{red}{what is that citation?}
\textcolor{red}{revisit: any quantitative way?}
